\chapter{Esercizi di analisi ammortizzata}\label{ch:appendice_esercizi_analisi_ammortizzata}

\section{Schema di costo risolutivo}\label{sec:appendice_esercizi_analisi_ammortizzata_schema}
Tutti gli esercizi girano attorno alla stessa forma. Dati tre parametri, $a > 0$ coefficiente di costo delle operazioni "costose", $b \ge 2$ intero, base delle operazioni costose, e $c > 0$ costo costante delle operazioni ordinarie, si definisce:
\begin{equation}
    c_i =
    \begin{cases}
        a \cdot i & \text{se } i \text{ è una potenza esatta di } b \\
        c & \text{altrimenti}
    \end{cases}
    \label{eq:appendice_esercizi_analisi_ammortizzata_schema_costo}
\end{equation}
da questo si chiede il costo ammortizzato per operazione.

Tutto gira attorno a una quantità $s$ che rappresenta la \textbf{quota che ogni operazione mette da parte per pagare la prossima operazione costosa}, ovvero è il costo delle potenze di $b$ spalmato uniformemente su tutte le operazioni. 

Questo deriva dal fatto che se prendiamo una potenza $i = b^k$, la sua operazione costa $ab^k$ (cioè molto). Questo costo viene pagato dalle operazioni che la precedono fino alla potenza precedente esclusa, più la potenza stessa, cioè dalle operazioni con indice in $(b^{k-1}, b^k]$\footnote{Ricordando che la parentesi tonda indica che l'estremo corrispondente è escluso dall'intervallo, mentre quella quadra indica che è incluso: qui $b^{k-1}$ è escluso e $b^k$ è incluso.}. Queste sono
\[
b^k - b^{k-1} = b^{k-1}(b-1)
\]

Se vogliamo che ognuna contribuisca con la stessa quota, dividiamo il costo da pagare per il numero di operazioni che se lo spartiscono:
\begin{equation}
s = \frac{ab^k}{b^{k-1}(b-1)} = \frac{ab}{b-1}
\label{eq:appendice_esercizi_analisi_ammortizzata_schema_quota}
\end{equation}

Poiché $k$ si semplifica, la quota non dipende dalla potenza che stiamo considerando: è la stessa per ogni blocco. Con tutti e tre i metodi (aggregazione, accantonamenti e potenziale) si ottiene che il costo ammortizzato per operazione è al più
\begin{equation}
    c + s = c + \frac{ab}{b-1}
    \label{eq:appendice_esercizi_analisi_ammortizzata_schema_costo_ammortizzato}
\end{equation}
cioè che $T(n)/n \le c+s$. Con gli accantonamenti e con il potenziale i singoli costi ammortizzati $\hat{c}_i$ non sono tutti uguali a $c+s$, ma nessuno lo supera.

\begin{osservazione}[Caso particolare della prima operazione]
    $i=1 = b^0$ è una potenza di $b$, quindi la prima operazione è costosa: il suo costo reale è $c_1 = a \cdot 1 = a$.
\end{osservazione}

\subsection{Metodo dell'aggregazione}\label{sec:appendice_esercizi_analisi_ammortizzata_aggregazione}

Ricordando la definizione \ref{def:aggregazione}, il metodo dell'aggregazione consiste nel calcolare il costo totale $T(n)$ di $n$ operazioni in sequenza e dividere per $n$. Se esprimiamo con $c_i$ il costo della $i$-esima operazione, il costo totale di $n$ operazioni è
\[
T(n) = \sum_{i=1}^n c_i = \sum_{\substack{i=1 \\ i \ \text{non potenza di } b}}^n c + \sum_{\substack{i=1 \\ i \ \text{potenza di } b}}^n a \cdot i
\]

\paragraph{1. Costo delle operazioni normali.}
Le operazioni normali costano $c$ ciascuna e sono \textbf{al più} $n$, perché in tutto le operazioni sono $n$\footnote{Questa è una maggiorazione larga, ma semplifica i conti.}. Quindi
\[
\sum_{i \ \text{non potenza di } b} c \leq cn
\]

\paragraph{2. Costo delle operazioni potenze.}
Al posto di contare le operazioni da $1$ a $n$, possiamo contare le potenze di $b$ a partire da $0$, ergo possiamo utilizzare $j$ come contatore. Questo significa che nella sommatoria
\[
\sum_{\substack{i=1 \\ i \ \text{potenza di } b}}^n a \cdot i
\]
compaiono solo gli $i$ della forma $i = b^j$. Ogni $i$ di questo tipo corrisponde a un \textbf{solo} esponente $j$, e viceversa. Possiamo quindi descrivere gli stessi termini usando $j$ al posto di $i$:
\[
i = b^j \implies a \cdot i = a \cdot b^j
\]
Quindi ci servono gli $i$ compresi tra $1$ e $n$, quindi gli $j$ tali che
\[
1 \le b^j \le n \implies \log_b 1 \le j \le \log_b n \implies 0 \le j \le \log_b n
\]

In generale $\log_b n$ non è un intero, quindi per contare gli $j$ validi dobbiamo prendere il floor di $\log_b n$. Quindi la sommatoria diventa
\[
\sum_{\substack{i=1 \\ i \ \text{potenza di } b}}^n a \cdot i = \sum_{j=0}^{\lfloor \log_b n \rfloor} a \cdot b^j
\]

Ricordando la formula della serie geometrica
\[
\sum_{j=0}^m b^j = \frac{b^{m+1}-1}{b-1}
\]
prendiamo $m = \lfloor \log_b n \rfloor$ e otteniamo
\[
\sum_{j=0}^{\lfloor \log_b n \rfloor} a \cdot b^j = a \cdot \frac{b^{\lfloor \log_b n \rfloor + 1}-1}{b-1}
\]

Consideriamo  due maggiorazioni:
\begin{enumerate}
    \item Buttiamo via il $-1$ al numeratore: sottraendo di meno la quantità può solo crescere, quindi otteniamo
    \[
    a \cdot \frac{b^{\lfloor \log_b n  \rfloor + 1} - 1}{b-1} \le a \cdot \frac{b^{\lfloor \log_b n \rfloor + 1}}{b-1}
    \]
    \item Spezziamo l'esponente in due parti, $b^{\lfloor \log_b n \rfloor + 1} = b^{\lfloor \log_b n \rfloor} \cdot b$. Per definizione di parte intera vale $\lfloor \log_b n \rfloor \le \log_b n$. Mettiamo entrambi i lati all'esponente di $b$: siccome $b > 1$, un esponente più grande dà un risultato più grande, quindi la disuguaglianza si conserva e, per definizione di logaritmo, $b^{\log_b n} = n$:
    \[
    \lfloor \log_b n \rfloor \le \log_b n \implies b^{\lfloor \log_b n \rfloor} \le b^{\log_b n} = n \implies b^{\lfloor \log_b n \rfloor + 1} \le b\,n
    \]
\end{enumerate}

Mettendo insieme le due maggiorazioni:
\[
a \cdot \frac{b^{\lfloor \log_b n \rfloor + 1} -1}{b-1} \le a \cdot \frac{b\,n}{b-1} = \frac{ab}{b-1} \cdot n = sn
\]

\paragraph{3. Aggregazione dei due costi.}
Sommando le due sommatorie otteniamo
\begin{equation}
\begin{aligned}
    T(n) =& \sum_{i=1}^n c_i \\
    =& \sum_{\substack{i=1 \\ i \ \text{non potenza di } b}}^n c + \sum_{\substack{i=1 \\ i \ \text{potenza di } b}}^n a \cdot i \\
    \le& cn + sn \\
    =& (c+s)n \\
    \implies& \frac{T(n)}{n} \le c+s
\end{aligned}
\label{eq:appendice_esercizi_analisi_ammortizzata_aggregazione_costo_totale}
\end{equation}

\begin{osservazione}[Il limite è stretto]
    Il limite $c+s$ non si può migliorare. Prendiamo $n = b^k$: le potenze di $b$ fino a $n$ sono le $k+1$ potenze $b^0, \dots, b^k$, quindi le operazioni ordinarie sono esattamente $n-k-1$ e
    \[
    T(n) = c\,(n-k-1) + a \cdot \frac{b^{k+1}-1}{b-1} = c\,(n-k-1) + a \cdot \frac{bn-1}{b-1}.
    \]
    Dividendo per $n$ si ha
    \[
    \frac{T(n)}{n} = c + s - \frac{c\,(k+1)}{n} - \frac{a}{(b-1)\,n} \xrightarrow{\ k \to \infty\ } c + s,
    \]
    perché $k = \log_b n$ cresce molto più lentamente di $n$.
\end{osservazione}

\subsection{Metodo degli accantonamenti}\label{sec:appendice_esercizi_analisi_ammortizzata_accantonamento}
Ricordando la definizione \ref{def:accantonamenti}, l'idea del metodo degli accantonamenti è scegliere dei costi ammortizzati $\hat{c}_i$ tali che il credito accumulato non diventi mai negativo. Da questo segue
\[
T(n) = \sum c_i \le \sum \hat{c}_i
\]

\paragraph{1. Scelta dei costi ammortizzati.}
Ogni operazione appartiene a un blocco $(b^{k-1}, b^k]$ e versa la quota $s$ per pagare la potenza che chiude il suo blocco. Ci sono tre casi:
\begin{description}
    \item[Operazione ordinaria] ($i$ non è potenza di $b$). Deve pagare il proprio costo reale $c$ e in più accantona la quota $s$:
    \[
    \hat{c}_i = c + s \qquad \text{se } i \text{ non è potenza di } b
    \]
    \item[Operazione costosa] ($i=b^k$ con $k \ge 1$). Anche lei fa parte del suo blocco, perché l'intervallo è chiuso a destra, quindi versa la quota $s$. Non deve pagare altro: il suo costo reale $ab^k$ è già stato coperto per intero dalle quote del blocco, compresa la sua. Quindi
    \[
    \hat{c}_i = s \qquad \text{se } i = b^k \text{ con } k \ge 1
    \]
    \item[Prima operazione] ($i=1=b^0$). Questa è un'operazione costosa, ma prima di lei non c'è nessuno che possa averla pagata. Non ha un blocco precedente, quindi paga da sola il proprio costo reale:
    \[
    \hat{c}_1 = c_1 = a
    \]
\end{description}

Da questo si estrapola
\begin{equation}
    \hat{c}_i =
    \begin{cases}
        a & \text{se } i=1 \\
        s & \text{se } i = b^k \text{ con } k \ge 1 \\
        c+s & \text{altrimenti}
    \end{cases}
\end{equation}

\paragraph{2. Il credito.}
Il credito dopo la $j$-esima operazione è
\[
B_j = \sum_{i=1}^j (\hat{c}_i - c_i)
\]

\noindent
Ogni operazione modifica il credito di $\hat{c}_i - c_i$:
\begin{itemize}
    \item un'operazione ordinaria lo \textbf{aumenta} di $(c+s)-c = s$,
    \item un'operazione costosa con $k \ge 1$ lo \textbf{diminuisce} di $ab^k -s$,
    \item la prima operazione non lo modifica, perché $a - a = 0$.
\end{itemize}

\noindent
Quindi bisogna dimostrare che $B_j \ge 0$ per ogni $j$. Per farlo dimostriamo per induzione su $k$ che, subito dopo l'operazione $b^k$, il credito vale $0$.
\begin{description}
    \item[Caso Base] ($k=0$). In questo caso $B_1 = \hat{c_1} - c_1 = a- a = 0$. Questo dimostra il caso base. 
    \item[Passo Induttivo] Sia $k \ge 1$ e supponiamo $B_{b^{k-1}} = 0$. Il blocco $(b^{k-1}, b^k]$ contiene $b^k - b^{k-1}$ operazioni, di cui $b^k - b^{k-1} - 1$ ordinarie e una costosa ($b^k$). Questo significa che il credito dopo le operazioni ordinarie, che aggiungono $s$ ciascuna, vale
    \[
    B_{b^k-1} = (b^k - b^{k-1}-1)s
    \]
    L'operazione costosa $b^k$ diminuisce il credito di $ab^k - s$, quindi il credito dopo l'operazione costosa vale
    \[
    B_{b^k} = (b^k - b^{k-1}-1)s - (ab^k - s) = (b^k - b^{k-1})s - ab^k
    \]
    Sostituendo $s = \frac{ab}{b-1}$ e $b^k - b^{k-1} = b^{k-1}(b-1)$ otteniamo
    \begin{align*}
        B_{b^k} &= s \cdot b^{k-1}(b-1) - ab^k \\
        &= \frac{ab}{b-1} \cdot b^{k-1}(b-1) - ab^k \\
        &= \frac{ab^k(b-1)}{b-1} - ab^k \\
        &= ab^k - ab^k \\
        &= 0
    \end{align*}
    Quindi il blocco raccoglie abbastanza credito dalle operazioni ordinarie per pagare l'operazione costosa, e il credito dopo l'operazione costosa vale $0$. Questo completa il passo induttivo.
\end{description}

Questo dimostra che il credito dopo ogni operazione costosa vale $0$. Tra una potenza e la successiva ci sono solo operazioni ordinarie, ognuna delle quali aggiunge $s$: se $b^k$ è l'ultima potenza di $b$ che non supera $j$, dopo l'operazione $b^k$ il credito vale $0$ e le operazioni $b^k+1, \dots, j$ ne aggiungono $s$ ciascuna, quindi
\[
B_j = s\,\big(j - b^{\lfloor \log_b j \rfloor}\big) \ge 0 \qquad \text{per ogni } j \ge 1.
\]
Quindi il credito è sempre non negativo.

\paragraph{3. Costo totale.}
Siccome $B_n \ge 0$, abbiamo che $\sum_{i=1}^n c_i \le \sum_{i=1}^n \hat{c}_i$.
Per ottenere un costo costante per operazione resta da maggiorare
$\sum_{i=1}^n \hat{c}_i$. I costi ammortizzati assumono solo tre valori,
$a$, $s$ e $c+s$, quindi basta maggiorare ognuno di essi con il più grande.
Questo è $c+s$, infatti:
\begin{itemize}
    \item $s \le c + s$, perché $c > 0$;
    \item $a \le c + s$, perché $s = a \cdot \frac{b}{b-1}$ e $\frac{b}{b-1} > 1$, quindi già $a < s$.
\end{itemize}
Quindi $\hat{c}_i \le c + s$ per ogni $i$ e otteniamo
\begin{equation}
    T(n) = \sum_{i=1}^n c_i \le \sum_{i=1}^n \hat{c}_i \le (c+s)\,n
    \implies \frac{T(n)}{n} \le c + s
\end{equation}
che è lo stesso costo ammortizzato trovato con il metodo dell'aggregazione.

\subsection{Metodo del Potenziale}\label{sec:appendice_esercizi_analisi_ammortizzata_potenziale}
Ricordando la definizione del metodo del potenziale \ref{def:potenziale}, l'idea è quella di associare a ogni configurazione della struttura un potenziale $\Phi_i$ tale che il costo ammortizzato della $i$-esima operazione sia
\[
    \hat{c}_i = c_i + \Phi_i - \Phi_{i-1}
\]

La somma dei costi ammortizzati è telescopica, quindi
\[
\sum_{i=1}^n \hat{c}_i = \sum_{i=1}^n c_i + \Phi_n - \Phi_0
\]
Se verifichiamo che $\Phi_n \ge \Phi_0$, otteniamo che $\sum_{i=1}^n \hat{c}_i \ge \sum_{i=1}^n c_i$, quindi il costo ammortizzato è una maggiorazione del costo reale. Perché questo valga per ogni prefisso della sequenza, e quindi per ogni $n$, basta dimostrare che $\Phi_i \ge \Phi_0$ per ogni $i$: se poniamo $\Phi_0 = 0$, la condizione diventa $\Phi_i \ge 0$.

\paragraph{1. Scelta del potenziale.}
Il problema principale nel metodo del potenziale è che non sappiamo a priori quale funzione usare come potenziale. Scegliamo allora il credito accumulato nel metodo degli accantonamenti: dopo la $j$-esima operazione, il potenziale è la somma delle quote $s$ versate dalle operazioni ordinarie eseguite dopo l'ultima potenza di $b$.

Sia $b^k$ l'ultima potenza di $b$ che non supera $j$ (indice del passo), cioè $b^k \le j < b^{k+1}$. Le operazioni ordinarie eseguite dopo di essa sono $b^k + 1, \dots, j$, quindi
$j - b^k$, e ognuna ha versato $s$:
\[
\Phi_j = s\,(j - b^k).
\]

Per esprimere $b^k$ in funzione di $j$, possiamo isolarlo applicando il logaritmo in base $b$:
\[
b^k \le j < b^{k+1} \implies k \le \log_b j < k+1
\]
questo dice che $k$ è un numero intero, non supera $\log_b j$ e il successivo $k+1$ lo supera già. Quindi
\[
k = \lfloor \log_b j \rfloor
\]

\noindent
Sostituendo in $\Phi_j = s(j - b^k)$ otteniamo:
\[
\Phi_j = 
\begin{cases}
    0 & \text{se } j = 0 \\
    s(j-b^{\lfloor \log_b j \rfloor}) & \text{se } j > 0
\end{cases}
\]

Il metodo funziona solo se $\Phi_j \ge \Phi_0 = 0, \ \forall j$. Per farlo riscriviamo
\[
s(j-b^{\lfloor \log_b j \rfloor}) \ge 0 \implies j - b^{\lfloor \log_b j \rfloor} \ge 0
\]
per la definizione di logaritmo, $b^{\log_b j} = j$, quindi $b^{\log_b j} - b^{\lfloor \log_b j \rfloor} \ge 0$. Poiché $b > 1$, la funzione $b^x$ è crescente, quindi $\log_b j - \lfloor \log_b j \rfloor \ge 0$, che è vera per definizione di parte intera. Quindi il potenziale è sempre non negativo:

\begin{osservazione}
    Non consideriamo il valore $s$ perché è una costante positiva ($a > 0$ e $b > 1$), quindi non può rendere negativo il potenziale nella sua moltiplicazione.
\end{osservazione}

\paragraph{2. Calcolo dei costi ammortizzati.}
I costi ammortizzati si calcolano come
\[
\hat{c}_i = c_i + \Phi_i - \Phi_{i-1}
\]
Il punto di questo calcolo rimane capire quanto vale $\Phi_i - \Phi_{i-1}$. Ci sono tre casi:
\begin{description}
    \item[Caso $i=1$] L'ultima potenza di $b$ che non supera $1$ è $b^0 = 1$, quindi $\Phi_1 = s(1-1) = 0$. Per $i=0$ abbiamo $\Phi_0 = 0$. Da questo segue che
    \[
    \hat{c}_1 = a + 0 - 0 = a
    \]
    \item[Caso $i$ non potenza di $b$] (ordinaria) Sia $b^k$ l'ultima potenza di $b$ che non supera $i$, cioè $b^k \le i < b^{k+1}$. Siccome $i$ non è una potenza, vale $b^k < i$ e quindi, essendo interi, $b^k \le i-1$. Inoltre $i-1 < i < b^{k+1}$. Allora anche $i-1$ sta in $[b^k, b^{k+1})$ e ha la stessa ultima potenza $b^k$:
    \begin{align*}
        \Phi_i - \Phi_{i-1} &= s(i-b^k) - s(i-1 - b^k) \\
        &= s(i-b^k - i + 1 + b^k) \\
        &= s
    \end{align*}
    quindi il costo ammortizzato vale
    \[
    \hat{c}_i = c + s
    \]
    Questo significa che un'operazione ordinaria paga il suo costo $c$ e aumenta il potenziale di una quota $s$.
    \item[Caso $i = b^k$ con $k \ge 1$] (costosa) Qui l'ultima potenza cambia:
    \begin{itemize}
        \item per $i$ l'ultima potenza è $i$ stessa, quindi $\Phi_i = s(i - b^{\lfloor \log_b i \rfloor}) = s(i - i) = 0$;
        \item per $i-1$ vale $b^{k-1} \le i-1 < b^k$: la disuguaglianza destra è ovvia, quella sinistra vale perché $b^{k-1} < b^k$ e, essendo interi, $b^{k-1} \le b^k - 1 = i-1$. Tra le due non ci sono altre potenze di $b$, perché sono consecutive. L'ultima potenza è quindi
        \[
        b^{k-1} = \frac{b^k}{b} = \frac{i}{b}
        \]
        e quindi
        \[
            \Phi_{i-1} = s\left(i-1 - \frac{i}{b}\right)
        \]
    \end{itemize}

    Poiché il costo reale di $c_i = ai$, il costo ammortizzato vale
    \begin{align*}
        \hat{c}_i &= c_i + \Phi_i - \Phi_{i-1} \\
        &= ai + 0 - s\left(i-1 - \frac{i}{b}\right) \\
        &= ai - s \cdot \frac{i(b-1)}{b} + s \\
        &= ai - \frac{ab}{b-1} \cdot \frac{i(b-1)}{b} + \frac{ab}{b-1} & \text{sostituendo } s = \frac{ab}{b-1} \\
        &= ai - ai \cdot \frac{b}{b} \cdot \frac{b-1}{b-1} + \frac{ab}{b-1} \\
        &= ai - ai + \frac{ab}{b-1} \\
        &= s & \text{risostituendo } \frac{ab}{b-1} = s
    \end{align*}
\end{description}

Ricapitolando, i costi ammortizzati sono:
\[
    \hat{c}_i = 
    \begin{cases}
        a & \text{se } i=1 \\
        s & \text{se } i = b^k \text{ con } k \ge 1 \\
        c+s & \text{altrimenti}
    \end{cases}
\]

\begin{osservazione}
    Da notare che i costi ammortizzati sono gli stessi del metodo degli accantonamenti, quindi il costo ammortizzato per operazione è lo stesso. Non è un caso: il potenziale scelto coincide esattamente con il credito, cioè $\Phi_j = B_j$ per ogni $j$.
\end{osservazione}

\paragraph{3. Costo totale.}
Poiché vale $\Phi_n \ge \Phi_0 = 0$, dalla somma telescopica segue $\sum c_i \le \sum \hat{c}_i$. Da qui si conclude, come negli accantonamenti, maggiorando ogni $\hat{c}_i$ con $c+s$ perché $a < s \le c+s$:
\begin{equation}
    T(n) = \sum_{i=1}^n c_i \le \sum_{i=1}^n \hat{c}_i \le (c+s)\,n
    \implies \frac{T(n)}{n} \le c + s
\end{equation}

\section{Esercizi}\label{sec:appendice_esercizi_analisi_ammortizzata_esercizi}

\begin{esercizio}\label{es:appendice_ammortizzata_potenze_5}
Utilizzando i tre metodi dell'analisi ammortizzata, si determini il costo ammortizzato per operazione di una sequenza di $n$ operazioni, ove il costo $c_i$ dell'$i$-esima operazione sia dato da
\[
c_i =
\begin{cases}
    12 \cdot i & \text{se } i \text{ è potenza esatta di } 5 \\
    7 & \text{altrimenti.}
\end{cases}
\]
\end{esercizio}

\subsubsection*{Soluzione}
L'esercizio ha la forma dello schema della sezione~\ref{sec:appendice_esercizi_analisi_ammortizzata_schema}, con
\[
a = 12, \qquad b = 5, \qquad c = 7.
\]

\paragraph{Aggregazione.}
Separiamo nella somma le operazioni ordinarie da quelle costose:
\[
T(n) = \sum_{i=1}^{n} c_i = \sum_{\substack{i=1 \\ i \ne 5^j}}^{n} 7 + \sum_{\substack{i=1 \\ i = 5^j}}^{n} 12 \cdot i .
\]
Le operazioni ordinarie sono al più $n$, quindi la prima somma vale al più $7n$. Nella seconda compaiono solo gli $i$ della forma $i = 5^j$, per i quali $12 \cdot i = 12 \cdot 5^j$; poiché $1 \le 5^j \le n$ equivale a $0 \le j \le \log_5 n$, cioè $0 \le j \le \lfloor \log_5 n \rfloor$, si ha
\[
\sum_{\substack{i=1 \\ i = 5^j}}^{n} 12 \cdot i = 12 \sum_{j=0}^{\lfloor \log_5 n \rfloor} 5^j = 12 \cdot \frac{5^{\lfloor \log_5 n \rfloor + 1} - 1}{5 - 1} \le 12 \cdot \frac{5n}{4} = 15n,
\]
dove abbiamo usato la somma geometrica e $5^{\lfloor \log_5 n \rfloor + 1} \le 5 \cdot 5^{\log_5 n} = 5n$. Sommando i due contributi otteniamo
\[
T(n) \le 7n + 15n = 22n \implies \frac{T(n)}{n} \le 22 .
\]

\paragraph{Accantonamenti.}
La potenza $5^k$ costa $12 \cdot 5^k$ e la pagano le $5^k - 5^{k-1} = 4 \cdot 5^{k-1}$ operazioni del blocco $(5^{k-1}, 5^k]$, quindi ognuna di esse mette da parte la quota
\[
s = \frac{12 \cdot 5^k}{4 \cdot 5^{k-1}} = 15 .
\]
I costi ammortizzati si scelgono distinguendo tre casi:
\begin{enumerate}
    \item $i = 1 = 5^0$: prima di lei nessuno può averla pagata, quindi paga da sola il proprio costo reale e non accantona niente, $\hat{c}_1 = c_1 = 12$;
    \item $i$ non potenza di $5$: paga il proprio costo reale e in più accantona la quota, $\hat{c}_i = 7 + 15 = 22$;
    \item $i = 5^k$ con $k \ge 1$: versa soltanto la quota $15$, perché il resto del suo costo è già stato accantonato dalle operazioni del suo blocco.
\end{enumerate}
Quindi
\[
\hat{c}_i =
\begin{cases}
    12 & \text{se } i = 1 \\
    15 & \text{se } i = 5^k \text{ con } k \ge 1 \\
    22 & \text{altrimenti.}
\end{cases}
\]
Resta da dimostrare che il credito $B_j = \sum_{i=1}^{j} (\hat{c}_i - c_i)$ non è mai negativo. Un'operazione ordinaria lo aumenta di $22 - 7 = 15$, mentre la potenza $5^k$ con $k \ge 1$ lo modifica di $15 - 12 \cdot 5^k < 0$: il credito può quindi scendere solo sulle potenze di $5$. Dimostriamo per induzione su $k$ che subito dopo ogni potenza il credito vale $0$, cioè $B_{5^k} = 0$.
\begin{description}
    \item[Caso base] ($k = 0$). $B_1 = \hat{c}_1 - c_1 = 12 - 12 = 0$.
    \item[Passo induttivo] Sia $k \ge 1$ e supponiamo $B_{5^{k-1}} = 0$. Il blocco $(5^{k-1}, 5^k]$ contiene $5^k - 5^{k-1} - 1$ operazioni ordinarie, che versano $15$ ciascuna, e l'operazione costosa $5^k$, quindi
    \begin{align*}
        B_{5^k} &= B_{5^{k-1}} + (5^k - 5^{k-1} - 1) \cdot 15 + (15 - 12 \cdot 5^k) \\
        &= (5^k - 5^{k-1}) \cdot 15 - 12 \cdot 5^k \\
        &= 5^{k-1}(5 - 1) \cdot 15 - 12 \cdot 5^k \\
        &= 60 \cdot 5^{k-1} - 12 \cdot 5 \cdot 5^{k-1} \\
        &= 60 \cdot 5^{k-1} - 60 \cdot 5^{k-1} = 0 .
    \end{align*}
\end{description}
Per i $j$ compresi tra due potenze, $5^{k-1} < j < 5^k$, dopo la potenza $5^{k-1}$ ci sono solo operazioni ordinarie, quindi
\[
B_j = B_{5^{k-1}} + 15\,(j - 5^{k-1}) = 15\,(j - 5^{k-1}) > 0 .
\]
Il credito è quindi sempre non negativo e vale $\sum_{i=1}^{n} c_i \le \sum_{i=1}^{n} \hat{c}_i$. Poiché $\hat{c}_i \le 22$ per ogni $i$,
\[
T(n) = \sum_{i=1}^{n} c_i \le \sum_{i=1}^{n} \hat{c}_i \le 22n \implies \frac{T(n)}{n} \le 22 .
\]

\paragraph{Potenziale.}
I costi ammortizzati sono $\hat{c}_i = c_i + \Phi_i - \Phi_{i-1}$. Come potenziale prendiamo il credito degli accantonamenti: se $5^k$ è l'ultima potenza di $5$ che non supera $j$, poniamo $\Phi_j = 15\,(j - 5^k)$. Per esprimere $5^k$ in funzione di $j$ isoliamo $k$:
\[
5^k \le j < 5^{k+1} \implies k \le \log_5 j < k + 1 \implies k = \lfloor \log_5 j \rfloor,
\]
da cui
\[
\Phi_j =
\begin{cases}
    0 & \text{se } j = 0 \\
    15\left(j - 5^{\lfloor \log_5 j \rfloor}\right) & \text{se } j > 0 .
\end{cases}
\]
Il potenziale non è mai inferiore a quello iniziale $\Phi_0 = 0$, infatti
\[
\lfloor \log_5 j \rfloor \le \log_5 j \implies 5^{\lfloor \log_5 j \rfloor} \le 5^{\log_5 j} = j \implies \Phi_j = 15\left(j - 5^{\lfloor \log_5 j \rfloor}\right) \ge 0 = \Phi_0,
\]
dove la seconda implicazione vale perché $5 > 1$. Calcoliamo i costi ammortizzati distinguendo tre casi:
\begin{description}
    \item[Caso $i = 1$] $\Phi_1 = 15\left(1 - 5^{\lfloor \log_5 1 \rfloor}\right) = 15\,(1 - 1) = 0$, quindi $\hat{c}_1 = 12 + 0 - 0 = 12$.
    \item[Caso $i$ non potenza di $5$] Sia $5^k$ l'ultima potenza di $5$ che non supera $i$. Poiché $i$ non è una potenza, $5^k \le i - 1$, quindi anche $i-1$ ha come ultima potenza $5^k$ e
    \[
    \Phi_i - \Phi_{i-1} = 15\,(i - 5^k) - 15\,(i - 1 - 5^k) = 15,
    \qquad
    \hat{c}_i = 7 + 15 = 22 .
    \]
    \item[Caso $i = 5^k$ con $k \ge 1$] L'ultima potenza che non supera $i$ è $i$ stessa, quindi $\Phi_i = 15\,(i - i) = 0$. Per $i - 1$ vale invece $5^{k-1} \le i - 1 < 5^k$, quindi la sua ultima potenza è quella precedente, $5^{k-1} = \frac{5^k}{5} = \frac{i}{5}$, e
    \[
    \Phi_{i-1} = 15\left(i - 1 - 5^{\lfloor \log_5 (i-1) \rfloor}\right) = 15\left(i - 1 - \frac{i}{5}\right).
    \]
    Allora
    \[
    \hat{c}_i = 12\,i + 0 - 15\left(i - 1 - \frac{i}{5}\right) = 12\,i - 15\,i + 15 + 3\,i = 15 .
    \]
\end{description}
Ritroviamo gli stessi costi ammortizzati degli accantonamenti:
\[
\hat{c}_i =
\begin{cases}
    12 & \text{se } i = 1 \\
    15 & \text{se } i = 5^k \text{ con } k \ge 1 \\
    22 & \text{altrimenti.}
\end{cases}
\]
Poiché $\Phi_n \ge \Phi_0$, dalla somma telescopica segue $\sum_{i=1}^{n} c_i \le \sum_{i=1}^{n} \hat{c}_i$, e maggiorando ogni $\hat{c}_i$ con il più grande dei tre valori otteniamo
\[
T(n) = \sum_{i=1}^{n} c_i \le \sum_{i=1}^{n} \hat{c}_i \le 22n \implies \frac{T(n)}{n} \le 22 .
\]

In sintesi, con tutti e tre i metodi il costo ammortizzato per operazione è al più $22 = c + s$, in accordo con lo schema generale.